% ============================================================================
% ASSIGNMENT TEMPLATE — ONE FILE, THREE MODES
%
% Change \AssignmentMode to one of:
%   submission : questions + solutions; student details; no rubrics
%   handout    : questions only; instructor/due date; optional answer space
%   key        : questions + solutions; optional grading rubrics
%
% Set \Gradedfalse to hide problem points and totals in every mode.
% Set \ShowTotalfalse to hide only the assignment total, keeping problem points.
% Rubrics appear ONLY when mode=key, \Gradedtrue, and \ShowRubricstrue.
% ============================================================================
\documentclass[11pt,a4paper]{article}

% ----------------------------------------------------------------------------
% 1. CONFIGURATION — edit these settings for each assignment
% ----------------------------------------------------------------------------
\providecommand{\AssignmentMode}{submission} % submission | handout | key; a \def before \input overrides it

\newif\ifGraded
\Gradedtrue                        % \Gradedtrue or \Gradedfalse
\newif\ifShowRubrics
\ShowRubricstrue                   % \ShowRubricstrue or \ShowRubricsfalse
\newif\ifShowTotal
\ShowTotaltrue                     % \ShowTotaltrue or \ShowTotalfalse

\newcommand{\CourseCode}{ECON 501}
\newcommand{\CourseName}{Microeconomic Theory}
\newcommand{\FullCourseName}{\CourseCode\ -- \CourseName}
\newcommand{\AcademicTerm}{Fall 2026}
\newcommand{\AssignmentType}{Problem Set}
\newcommand{\AssignmentNumber}{1}
\newcommand{\AssignmentTopic}{Optimization and Convexity}
\newcommand{\InstructorName}{Instructor Name}
\newcommand{\StudentName}{Student Name}
\newcommand{\StudentID}{123456789}
\newcommand{\DueDate}{9 September 2026}
\newcommand{\SubmissionDate}{9 October 2026}
\newcommand{\AssignmentInstructions}{Show your reasoning and state any assumptions you use.}

% ----------------------------------------------------------------------------
% 2. PACKAGES AND FORMATTING — normally leave unchanged
% ----------------------------------------------------------------------------
\usepackage[T1]{fontenc}
\usepackage{lmodern}
\usepackage{amsmath,amssymb,mathtools}
\usepackage[margin=0.85in]{geometry}
\usepackage{enumitem}
\usepackage{etoolbox}
\usepackage{fancyhdr}
\usepackage[hidelinks]{hyperref}

\setlength{\parindent}{0pt}
\setlength{\parskip}{0.32em}
\setlength{\headheight}{14pt}
\setlist[enumerate,1]{label = (\alph*),leftmargin=*,itemsep=0.35em}
\setlist[enumerate,2]{label = (\roman*),leftmargin=*,itemsep=0.2em}

% ----------------------------------------------------------------------------
% 3. MODE LOGIC — do not edit for ordinary use
% ----------------------------------------------------------------------------
\newif\ifSubmissionMode
\newif\ifHandoutMode
\newif\ifKeyMode
\newif\ifShowSolutions
\newif\ifPrintRubrics

\ifdefstring{\AssignmentMode}{submission}{%
    \SubmissionModetrue
    \ShowSolutionstrue
    \newcommand{\ModeLabel}{Submission}
}{%
    \ifdefstring{\AssignmentMode}{handout}{%
        \HandoutModetrue
        \newcommand{\ModeLabel}{Handout}
    }{%
        \ifdefstring{\AssignmentMode}{key}{%
            \KeyModetrue
            \ShowSolutionstrue
            \newcommand{\ModeLabel}{Solution Key}
        }{%
            \PackageError{assignment}
            {Invalid AssignmentMode: \AssignmentMode}
            {Use submission, handout, or key.}%
        }%
    }%
}

\ifKeyMode
\ifGraded
\ifShowRubrics
\PrintRubricstrue
\fi
\fi
\fi

% ----------------------------------------------------------------------------
% 4. HEADER
% ----------------------------------------------------------------------------
\pagestyle{fancy}
\fancyhf{}
\fancyhead[L]{\small\FullCourseName}
\fancyhead[R]{\small\AssignmentType\ \AssignmentNumber\enspace|\enspace\ModeLabel}
\fancyfoot[C]{\small\thepage}
\renewcommand{\headrulewidth}{0.4pt}
\fancypagestyle{plain}{%
    \fancyhf{}
    \fancyfoot[C]{\small\thepage}
    \renewcommand{\headrulewidth}{0pt}%
}

\newcommand{\MakeAssignmentHeader}{%
    \thispagestyle{plain}%
    {\small\bfseries\FullCourseName}%
    \ifdefempty{\AcademicTerm}{}{\hfill{\small\AcademicTerm}}%
    \par\vspace{0.45em}%
    {\LARGE\bfseries\AssignmentType}%
    \ifdefempty{\AssignmentNumber}{}{\ {\LARGE\bfseries\AssignmentNumber}}%
    \ifKeyMode\enspace{\large\normalfont(Solution Key)}\fi
    \par
    \ifdefempty{\AssignmentTopic}{}{%
        \vspace{0.2em}{\large\AssignmentTopic}\par
    }%
    \vspace{0.40em}%
    {\small
        \ifSubmissionMode
        \ifdefempty{\StudentName}{}{\textbf{Name:} \StudentName\qquad}%
        \ifdefempty{\StudentID}{}{\textbf{ID:} \StudentID\qquad}%
        \ifdefempty{\SubmissionDate}{}{\textbf{Submitted:} \SubmissionDate}%
        \else
        \ifdefempty{\InstructorName}{}{\textbf{Instructor:} \InstructorName\qquad}%
        \ifdefempty{\DueDate}{}{\textbf{Due:} \DueDate}%
        \fi
    }\par
    \vspace{0.45em}\hrule\vspace{0.40em}%
    \ifdefempty{\AssignmentInstructions}{}{%
        {\small\textbf{Instructions.} \AssignmentInstructions\par}%
        \vspace{0.35em}%
    }%
}

% ----------------------------------------------------------------------------
% 5. PROBLEM ENVIRONMENTS
% ----------------------------------------------------------------------------
\newcounter{problem}
\newcounter{totalpoints}

% Usage: \begin{problem}[points][optional title] ... \end{problem}
% The point count is a nonnegative integer; the default is zero.
\NewDocumentEnvironment{problem}{O{0} O{}}{%
    \par\addvspace{0.85\baselineskip}%
    \refstepcounter{problem}%
    \ifnum#1<0\relax
    \PackageError{assignment-standalone}{Problem points cannot be negative}%
    {Provide a nonnegative integer for points.}%
    \fi
    \addtocounter{totalpoints}{#1}%
    \noindent\textbf{Problem~\theproblem}%
    \ifstrempty{#2}{}{\textbf{ --- #2}}%
    \ifGraded
    \ifnum#1>0\relax\hfill{\small\textbf{(#1 points)}}\fi
    \fi
    \par\nobreak\smallskip\ignorespaces
}{%
    \par\addvspace{0.5\baselineskip}%
}

% Body collection lets hidden solutions and rubrics remain in the same source.
\NewDocumentEnvironment{solution}{+b}{%
    \ifShowSolutions
    \par\smallskip\noindent\textbf{Solution.}\par\nobreak\smallskip
    #1\par
    \fi
}{}

\NewDocumentEnvironment{rubric}{+b}{%
    \ifPrintRubrics
    \par\smallskip\noindent\textbf{Grading rubric.}\par\nobreak\smallskip
    {\small #1\par}%
    \fi
}{}

% Optional writing space, used only in handout mode.
\NewDocumentCommand{\answerspace}{m}{%
    \ifHandoutMode\par\vspace{#1}\fi
}

% Place after the final problem. Displayed only if graded and ShowTotal is true.
\NewDocumentCommand{\printtotal}{}{%
    \ifGraded
    \ifShowTotal
    \par\bigskip\hrule\smallskip
    \noindent{\small\bfseries Total: \arabic{totalpoints} points}\par
    \fi
    \fi
}

% ============================================================================
% 6. ASSIGNMENT CONTENT — write each question just once
% ============================================================================
\begin{document}
\MakeAssignmentHeader

\begin{problem}\label{prob:consumer}
    A consumer chooses \(x,y>0\) to maximize
    \[
        u(x,y)=\log x+\log y
    \]
    subject to \(p_x x+p_y y=m\), where \(p_x,p_y,m>0\).

    \begin{enumerate}
        \item Write the Lagrangian and the first-order conditions.
        \item Find the optimal consumption bundle \((x^*,y^*)\).
    \end{enumerate}

    \answerspace{3cm} % Blank working space in handout mode only.

    \begin{solution}
        Solution
    \end{solution}

    \begin{rubric}
        \begin{itemize}[leftmargin=*,nosep]
            \item Lagrangian and first-order conditions: 4 points.
            \item Correct demand functions: 4 points.
            \item Justification of optimality: 2 points.
        \end{itemize}
    \end{rubric}
\end{problem}

\begin{problem}[5][Convex combinations]
    Let \(f\) and \(g\) be convex functions on a convex set
    \(C\subseteq\mathbb{R}^n\). Prove that
    \(h=\alpha f+(1-\alpha)g\) is convex for every \(\alpha\in[0,1]\).

    \answerspace{2cm}

    \begin{solution}
        For any \(x,y\in C\) and \(t\in[0,1]\), convexity and nonnegative
        weights imply
        \begin{align*}
            h(tx+(1-t)y)
            &\leq \alpha\bigl(tf(x)+(1-t)f(y)\bigr)
            +(1-\alpha)\bigl(tg(x)+(1-t)g(y)\bigr)\\
            &=t h(x)+(1-t)h(y).
        \end{align*}
        Hence \(h\) is convex.
    \end{solution}

    \begin{rubric}
        \begin{itemize}[leftmargin=*,nosep]
            \item Applies convexity of both \(f\) and \(g\): 3 points.
            \item Correctly combines the inequalities: 2 points.
        \end{itemize}
    \end{rubric}
\end{problem}

\begin{problem}[10][Consumer choice]\label{prob:consumer}
    A consumer chooses \(x,y>0\) to maximize
    \[
        u(x,y)=\log x+\log y
    \]
    subject to \(p_x x+p_y y=m\), where \(p_x,p_y,m>0\).

    \begin{enumerate}
        \item Write the Lagrangian and the first-order conditions.
        \item Find the optimal consumption bundle \((x^*,y^*)\).
    \end{enumerate}

    \answerspace{3cm} % Blank working space in handout mode only.

    \begin{solution}
        Solution
    \end{solution}

    \begin{rubric}
        \begin{itemize}[leftmargin=*,nosep]
            \item Lagrangian and first-order conditions: 4 points.
            \item Correct demand functions: 4 points.
            \item Justification of optimality: 2 points.
        \end{itemize}
    \end{rubric}
\end{problem}

\begin{problem}[5][Convex combinations]
    Let \(f\) and \(g\) be convex functions on a convex set
    \(C\subseteq\mathbb{R}^n\). Prove that
    \(h=\alpha f+(1-\alpha)g\) is convex for every \(\alpha\in[0,1]\).

    \answerspace{2cm}

    \begin{solution}
        For any \(x,y\in C\) and \(t\in[0,1]\), convexity and nonnegative
        weights imply
        \begin{align*}
            h(tx+(1-t)y)
            &\leq \alpha\bigl(tf(x)+(1-t)f(y)\bigr)
            +(1-\alpha)\bigl(tg(x)+(1-t)g(y)\bigr)\\
            &=t h(x)+(1-t)h(y).
        \end{align*}
        Hence \(h\) is convex.
    \end{solution}

    \begin{rubric}
        \begin{itemize}[leftmargin=*,nosep]
            \item Applies convexity of both \(f\) and \(g\): 3 points.
            \item Correctly combines the inequalities: 2 points.
        \end{itemize}
    \end{rubric}
\end{problem}

%\printtotal
\end{document}